\begin{proofbox}
  \textbf{\textcolor{CProof}{证明前先记一句话}}\par\nopagebreak[4]

  对任何可导函数，\textbf{两边相等的式子可以同时求导}；反过来，若某个函数的导数恒为 $0$，它就是一个常数。下面三种证明都在高中导数知识内。

  \medskip
  \textbf{\textcolor{CProof}{方法一（最直接）：原函数轴对称 $\Longrightarrow$ 导函数中心对称}}\par\nopagebreak[4]

  \textbf{思路：}把左右对称的横坐标写成 $a+t$ 和 $a-t$，两边对 $t$ 求导。
  \begin{description}[style=nextline,leftmargin=2.8em,labelsep=0.5em,itemsep=0.35em]
    \item[\textbf{第一步（写出“同高”）：}] 原函数关于 $x=a$ 对称，所以 $f(a+t)=f(a-t)$。
    \item[\textbf{第二步（分别求导）：}] 左边对 $t$ 求导是 $f'(a+t)$；右边因为里面是 $a-t$，还要乘 $-1$，得到 $-f'(a-t)$。
    \item[\textbf{第三步（读出对称）：}] $f'(a+t)=-f'(a-t)$，左右导数正负相反，导函数关于 $(a,0)$ 中心对称。
  \end{description}
  \[
    \boxed{f(a+t)=f(a-t)\ \Longrightarrow\ f'(a+t)=-f'(a-t).}
  \]
  \begin{center}
  \begin{tikzpicture}[font=\small,>=stealth,line join=round]
\draw[draw=teal!45!black,fill=green!3,rounded corners=3pt] (0,0) rectangle (13.6,6.1);
\node at (6.8,5.55) {$a-t$ 中的负号：$t$ 增大时，这个输入向左移动};
\begin{scope}[shift={(1.3,1.6)},x=1.2cm,y=0.65cm]
\draw[gray,->] (-0.1,0)--(4.4,0) node[right] {$x$};
\draw[gray,->] (0,-0.2)--(0,4.8);
\draw[teal!75!black,very thick,domain=0.3:3.7,samples=100] plot(\x,{(\x-2)^2+1});
\fill[teal!75!black] (1,2) circle(2.5pt);
\fill[teal!75!black] (3,2) circle(2.5pt);
\draw[dashed,gray] (2,0)--(2,4.5);
\draw[->,orange!85!black,very thick] (1,-.8)--(.55,-.8);
\draw[->,blue!70!black,very thick] (3,-.8)--(3.45,-.8);
\node[below] at (1,0) {$a-t=1$};\node[below] at (3,0) {$a+t=3$};
\end{scope}
\node[anchor=west,align=left,text width=5.4cm] at (7.3,3) {取 $a=2,t=1$。\\[10pt]$t$ 增大时：\\左点向左，右点向右；\\两边都沿图像升高。\\[10pt]但 $f^{\prime}(1)=-2$、$f^{\prime}(3)=2$。\\负号来自左边输入的反向变化。};
\node at (6.8,.4) {橙箭头：$a-t$ 减小；蓝箭头：$a+t$ 增大。};
\end{tikzpicture}

  \end{center}

  \medskip
  \textbf{\textcolor{CProof}{方法二（反向也成立）：利用“导数为零”的常数规律}}\par\nopagebreak[4]

  \textbf{A．导函数中心对称，反推原函数轴对称。}已知 $f'(a+t)=-f'(a-t)$。令
  \[
    H(t)=f(a+t)-f(a-t).
  \]
  则 $H'(t)=f'(a+t)+f'(a-t)=0$，所以 $H(t)$ 是常数。由于 $H(0)=f(a)-f(a)=0$，因此 $H(t)=0$，即 $f(a+t)=f(a-t)$。

  \textbf{B．导函数轴对称，反推原函数中心对称。}已知 $f'(a+t)=f'(a-t)$。令
  \[
    K(t)=f(a+t)+f(a-t).
  \]
  则 $K'(t)=f'(a+t)-f'(a-t)=0$，所以 $K(t)$ 是常数。代入 $t=0$，$K(0)=2f(a)$，得到
  \[
    \boxed{f(a+t)+f(a-t)=2f(a).}
  \]
  同样，从这个等式两边对 $t$ 求导又可推出 $f'(a+t)=f'(a-t)$，因此第二组也是\textbf{双向成立}。
  \begin{center}
  \begin{tikzpicture}[font=\small,>=stealth,line join=round]
\draw[draw=teal!45!black,fill=green!3,rounded corners=3pt] (0,0) rectangle (13.6,5.9);
\node at (6.8,5.3500000000000005) {反推的关键：$H$、$K$ 的导数为零，各自图像都是水平线};
\begin{scope}[shift={(4.8,1.3)},x=1.35cm,y=0.18cm]
\draw[gray,->] (-2.6,0)--(2.7,0) node[right] {$t$};
\draw[gray,->] (0,-1)--(0,17);
\draw[teal!75!black,very thick] (-2.3,0)--(2.3,0);
\draw[blue!70!black,very thick] (-2.3,14)--(2.3,14);
\draw[gray,dashed] (0,0)--(0,14);
\fill[teal!75!black] (0,0) circle(2.5pt);
\fill[blue!70!black] (0,14) circle(2.5pt);
\node[below] at (0,-.2) {$0$};
\end{scope}
\node[anchor=west,align=left,text width=3.8cm] at (9,3.8) {$K(t)\equiv14$\\取 $f(x)=(x-1)^3+7$\\$a=1,\ K(0)=2f(1)=14$。};
\node[anchor=west,align=left,text width=3.8cm] at (9,1.6) {$H(t)\equiv0$\\取 $f(x)=(x-2)^2+1$\\$a=2,\ H(0)=0$。};
\node at (6.8,.4) {图示用具体函数展示；正文证明对任意符合条件的 $f$ 成立。};
\end{tikzpicture}

  \end{center}

  \medskip
  \textbf{\textcolor{CProof}{方法三（完整推广）：导函数中心 $(a,b)$ 的一般情形}}\par\nopagebreak[4]

  \textbf{第一步：}导函数关于 $(a,b)$ 中心对称，按中心对称的定义，
  \[
    f'(a+t)+f'(a-t)=2b.
  \]
  \textbf{第二步：}令 $P(t)=f(a+t)-f(a-t)-2bt$，求导得
  \[
    P'(t)=f'(a+t)+f'(a-t)-2b=0.
  \]
  \textbf{第三步：}既然 $P(t)$ 是常数，再代入 $t=0$，$P(0)=0$，就得到
  \[
    \boxed{f(a+t)-f(a-t)=2bt.}
  \]
  反过来，对上式求导，又能得到原来的导函数中心对称关系。特别是 $b=0$，立刻回到第一组规律。
  \begin{center}
  \begin{tikzpicture}[font=\small,>=stealth,line join=round]
\draw[draw=teal!45!black,fill=green!3,rounded corners=3pt] (0,0) rectangle (13.6,6.5);
\node at (6.8,5.95) {导数的平均数是 $b$，原函数对应两点的高度差是 $2bt$};
\node at (3.5,5.2) {$f^{\prime}(x)=2x+2$};
\node at (10.2,5.2) {$f(x)=x^2+2x$};
\begin{scope}[shift={(3.3,1.75)},x=1.25cm,y=0.6cm]
\draw[gray,->] (-1.8,0)--(1.8,0) node[right] {$x$};
\draw[gray,->] (0,-1.2)--(0,4.6);
\draw[blue!70!black,very thick,domain=-1.45:1.2,samples=100] plot(\x,{2*\x+2});
\fill[blue!70!black] (-1,0) circle(2.5pt);
\fill[blue!70!black] (1,4) circle(2.5pt);
\fill[blue!70!black] (0,2) circle(2.5pt);
\draw[dotted,gray] (-1,0)--(1,4);
\node[left] at (-.1,2) {$(0,2)$};
\node[above left] at (-1,.1) {$0$};\node[above right] at (1,4.1) {$4$};
\draw[gray] (-1,0+.08)--(-1,0-.08);\node[below] at (-1,0-.15) {$-1$};
\draw[gray] (1,0+.08)--(1,0-.08);\node[below] at (1,0-.15) {$1$};
\end{scope}
\begin{scope}[shift={(9.4,2.2)},x=1.25cm,y=0.6cm]
\draw[gray,->] (-1.8,0)--(2,0) node[right] {$x$};
\draw[gray,->] (0,-1.7)--(0,4.1);
\draw[teal!75!black,very thick,domain=-1.65:1.1,samples=100] plot(\x,{(\x)^2+2*\x});
\fill[teal!75!black] (-1,-1) circle(2.5pt);
\fill[teal!75!black] (1,3) circle(2.5pt);
\draw[orange!85!black,<->,thick] (1.55,-1)--(1.55,3);
\node[right] at (1.6,1) {$4$};
\draw[gray,dotted] (-1,-1)--(1.55,-1) (1,3)--(1.55,3);
\node[below left] at (-1,-1.1) {$-1$};\node[above] at (1,3.25) {$3$};
\draw[gray] (-1,0+.08)--(-1,0-.08);\node[below] at (-1,0-.15) {$-1$};
\draw[gray] (1,0+.08)--(1,0-.08);\node[below] at (1,0-.15) {$1$};
\end{scope}
\node at (6.8,.4) {取 $a=0,b=2,t=1$：$f(1)-f(-1)=3-(-1)=4=2bt$。};
\end{tikzpicture}

  \end{center}
  \textbf{结论回扣：}两种对称关系能互推，靠的就是“\textbf{求导时左右差一个负号}”和“\textbf{导数恒为零则原式为常数}”。
\end{proofbox}
